% Generated draft: requires mathematical and source review.
\documentclass[11pt]{article}
\usepackage[margin=25mm]{geometry}
\usepackage{fontspec}
\setmainfont{Latin Modern Roman}
\usepackage{amsmath,amssymb,mathtools,unicode-math}
\setmathfont{Latin Modern Math}
\usepackage{enumitem,needspace,xurl,hyperref}
\setlength{\parindent}{0pt}
\setlength{\parskip}{5pt}
\setlength{\emergencystretch}{3em}
\newcommand{\N}{\mathbb{N}}
\newcommand{\Z}{\mathbb{Z}}
\newcommand{\Q}{\mathbb{Q}}
\newcommand{\R}{\mathbb{R}}
\newcommand{\C}{\mathbb{C}}
\newcommand{\F}{\mathbb{F}}
\newcommand{\A}{\mathbb{A}}
\newcommand{\PP}{\mathbb{P}}
\newcommand{\E}{\mathbb{E}}
\newcommand{\eps}{\varepsilon}
\newcommand{\dd}{\,\mathrm{d}}
\begin{document}
\section*{UnsolvedMath ID: 701551}
\textbf{Source problem code:} catalogue entry
\subsection{Definitions and mathematical statement}
Let \(A\) and \(B\) be groups with specified embeddings of a common group \(H\), and let \(G=A\ast_H B\) be their free product with amalgamation. Identify \(A\), \(B\), and \(H\) with their canonical images in \(G\). Write \(1\) for the trivial subgroup and \(|A:H|\), \(|B:H|\) for subgroup indices. The minimum condition on subgroups of \(H\) means that every descending chain \(H_1\ge H_2\ge\cdots\) of subgroups of \(H\) is eventually constant. Saying that \(H\) is a retract of both factors means that there are homomorphisms \(r_A:A\to H\) and \(r_B:B\to H\) whose restrictions to \(H\) are the identity.

\textbf{Unresolved source notation.} The supplied record does not define \(\psi\), \(\lambda\), \(H_G\), or the phrase ``nearly splits over \(H\).'' These are retained below as source notation without assigning them meanings. Consequently, the following faithfully records the proposed implications but is not yet a fully specified mathematical statement. The record also imposes no explicit requirement that either embedding of \(H\) into a factor be proper.

For each of the following seven conditions, separately, the question asks whether every amalgamated free product as above satisfying that condition must satisfy \(\psi(G)=1\):
\begin{enumerate}
\item[(a)] \(H\) is finite cyclic and \(H_G=1\).
\item[(b)] \(H\) is finite cyclic and
\[
\lambda(A)\cap H_G=1\quad\text{or}\quad\lambda(B)\cap H_G=1.
\]
\item[(c)] \(H_G=1\) and \(H\) satisfies the minimum condition on subgroups.
\item[(d)] \(\lambda(G)\cap H=1\).
\item[(e)] \(A\) and \(B\) are free groups, \(H\) is finitely generated, and at least one of \(|A:H|\) and \(|B:H|\) is infinite.
\item[(f)] \(H\) is infinite cyclic and is a retract of both \(A\) and \(B\).
\item[(g)] \(G\) nearly splits over \(H\), and \(H\trianglelefteq G\) has order \(p\) for some prime \(p\).
\end{enumerate}
Each condition is an alternative hypothesis for the same question, not an additional hypothesis to be combined with the other six. The undefined notation and splitting terminology require clarification from the cited source before these questions can be made fully precise.
\subsection{Short English statement}
M. K. Azarian asks whether \(\psi(G)\) is trivial for an amalgamated free product \(G=A\ast_H B\) under each of seven separate hypotheses on the factors and amalgamated subgroup. The supplied record leaves \(\psi\), \(\lambda\), \(H_G\), and ``nearly splits'' undefined, so the intended questions require source clarification. The registry labels the problem open; that label is supplied metadata, not an independently verified assessment.
\subsection{Sources}
\begin{itemize}
\item Local UnsolvedMath registry, supplied complete source record: permanent ID \(701551\), exact source problem code \texttt{catalogue entry}, title ``Kourovka Notebook Problem 19.3.'' Upstream revision: \texttt{d371ee917ac9647fbb06cd6fc529deec6f3dda82}. The supplied status is \texttt{open}; it has not been independently verified.
\item The supplied record attributes the problem to M. K. Azarian, proposed in 2018, and cites \emph{The Kourovka Notebook (Unsolved Problems in Group Theory)}, 19th issue (2018), original problem 19.3.
\item The supplied background additionally cites the 21st issue edition (2026), \texttt{arXiv:1401.0300v48}. The supplied collection description names E. I. Khukhro and V. D. Mazurov as editors. These bibliographic details are supplied attribution; the cited editions have not been inspected for this draft.
\end{itemize}
\end{document}
